\documentclass{article}
\usepackage{graphicx} % Required for inserting images
\usepackage{graphicx, amsmath, amssymb, siunitx}
\usepackage[left=4cm, right=4cm]{geometry}
\usepackage{listings}
\usepackage{xcolor}
\usepackage{url}
% Custom styling for Python code
\lstset{
    language=Python,
    basicstyle=\ttfamily\small,
    keywordstyle=\color{blue}\bfseries,
    stringstyle=\color{red},
    commentstyle=\color{green!50!black}\itshape,
    numbers=left,
    numberstyle=\tiny\color{gray},
    stepnumber=1,
    breaklines=true,
    frame=single,
    showstringspaces=false
}

\title{Building a Neural Network from Scratch}
\author{Pratham Kelkar}
\date{July 2026}

\begin{document}

\maketitle

\pagebreak

\tableofcontents
\pagebreak

\section{Motivation}

Neural networks form the backbone of many technical applications we use today from FaceID on unlocking our phones to talking to artificial intelligence assistants like ChatGPT and Gemini. 
\\
\\
I have always been fascinated by how they work under the hood. Before this project, I knew that neural networks take in an input and perform some sort of multi-variable regression or classification on the input to provide an output. Although I understood the underlying mechanism that they work on, I did not know the rigorous mathematical concepts that they operate on and how one is actually implemented without the help of frameworks like TensorFlow and PyTorch. 
\\
\\
Previously, I have experimented with these frameworks and created accurate models by piling on dense and dropout layers, but I never understood the underlying concepts or the implementation details. This is what I aim to develop through this project, hoping that it will leave me in a better position to create and optimize models in the future.
\\
\\
For this project, I will be using the ATLAS Higgs Boson Challenge dataset, originally released by CERN and hosted through a Kaggle competition (``Higgs Boson Machine Learning Challenge''), to train and test my models. This makes it a good fit for training and testing, since it has $250,000$ rows of data spanning $30$ features, along with the fact that it has been widely used and benchmarked before, giving me a good point of reference to compare my own results against.

\section{Introduction to Neural Networks}
Neural networks are essentially machine learning algorithms used to predict the output of a given input by performing mathematical operations on it. Mathematically, a neural network learns a function $f(\mathbf{x})$ by mapping an input vector $\mathbf{x} = [x_1, x_2, x_3, \dots, x_n]$ to predict a response $\mathbf{y}$. 

The simplest form of a neural network is a linear regression model, which can be thought of as a neural network with a single neuron in it. Essentially, all neurons of a neural network perform a linear transformation on the input variable to provide an output that minimizes the overall loss of the predictions it makes.

\subsection{The Problem to be Solved}
Many neural networks try to solve classification problems (for example: Google Image Search or FaceID on a phone). They process inputs and "classify" them into one of many (or in some cases just two) bins. I will be focusing on the framework for binary classification problems, which can then be modified to fit multi-variable classification. 

\subsection{Binary Classification}
The binary classification problem focuses on classifying an input $\mathbf{x}$ as a $1$ (True) or a $0$ (False) as a prediction $\hat{y}$ made by the model. For example, if we take up the problem of identifying whether an image contains a cat, the model will output a 1 for images that it predicts contains a cat and a 0 for images that it predicts to not contain a cat. It processes an input with a set of features to predict the probability that the given input is True or False. If this probability exceeds a numerical threshold $\epsilon$, then it can be classified as a 1.\\
In mathematical terms, given an input $\mathbf{x}$, we find
$$\hat{y} = P(y=1|\mathbf{x})$$
The simplest way to build such a neural network is by understanding the logistic regression algorithm as a base for the process that each neuron executes.

\section{Logistic Regression}

\subsection{Mathematical Representation}

Logistic regression is a simple algorithm that applies a linear transformation to an input and then converts it into a probability using the sigmoid function. It repeats this process for multiple iterations to minimize the error or the "loss" from each prediction.
\\
Taking $\mathbf{x} = [x_1, x_2, \dots, x_n], \mathbf{x} \in \mathbb{R}^n$ as an input vector with a set of features, we want to find the output $\hat{y}$ such that:
$$\hat{y} = P(y=1|\mathbf{x})$$
Here, we apply a linear transformation of the form $y=mx+c$ to $\mathbf{x}$, using the parameters $\mathbf{w}=[w_1, w_2, \dots, w_n]$ and $b$, such that $\mathbf{w} \in \mathbb{R}^n$ and $b \in \mathbb{R}$. This can be represented as:
\begin{equation}
    z = \mathbf{w}^T\mathbf{x}+b
\end{equation}
\subsubsection{Activation Functions}
Now that we have the linear part $z$, we need to turn it into a prediction. On its own, the value $z = \mathbf{w}^T\mathbf{x} + b$ is unbounded and can take any value in $\mathbb{R}$, so it cannot be read directly as a probability. To fix this, we pass $z$ through a non-linear function called an ``activation function'', which is written as $g(\cdot)$, so that the output of a neuron is $a = g(z)$.
\\
\\
An activation function serves two purposes. First, it maps the unbounded output of the linear step into a useful range. Second, and more importantly for the more complicated networks with multiple layers, it introduces non-linearity into the model: if we only composed linear transformations, then stacking many neurons on top of one another would still collapse into a single linear transformation, regardless of how many were used, and the model could never capture the complex, curved patterns that real data demands. The activation function is what breaks this linearity.
For logistic regression, the activation we choose is the sigmoid function, defined as:
\begin{equation}
    a = g(z) = \sigma(z) = \frac{1}{1 + e^{-z}}
\end{equation}
The sigmoid squashes any real input into the open interval $(0, 1)$, which is exactly what we need in order to interpret the output as the probability that the input belongs to the positive class. Thus,
$$a = \sigma(z) = \hat{y}$$
\subsubsection{Loss and Cost Function}
The error in the prediction can be represented by:
$$y - \hat{y}$$
Thus, we can now define a cost function that we can minimize to ensure that out predictions keep getting more accurate with each iteration of this process.
\\
For the loss function, the squared error function $(\hat{y}-y)^2$ is not used because it leads to a non-convex optimization problem (being a function of $\mathbf{w}$ and b) with multiple optima where the algorithm could get stuck, and thus the global minima may not be found (Ng). Instead, we can define the loss function as follows:
\begin{equation}
    \mathcal{L}(\hat{y},y)=-(y\log \hat{y} + (1-y)\log{(1-\hat{y}))}
\end{equation}
Here, if $y=1$, $\mathcal{L}(\hat{y},y)=-\log \hat{y}$ and if $y=0$, $\mathcal{L}(\hat{y},y)=-\log (1-\hat{y})$.
\\Now, taking $\hat{y}=a$, the cost function for $1$ training example can be defined as:
\begin{equation}
    J(\mathbf{w},b)=\mathcal{L}(\hat{y},y) = -[y\log a + (1-y)\log{(1-a)}]
\end{equation}
Now that the cost function $J(\mathbf{w}, b)$ is defined, it needs to be minimized to train the model on the set of $m$ training examples so that it can accurately classify them. This can be done by taking the gradient descent step. This is a small step to update the parameters $\mathbf{w}$ and $b$ in the direction of steepest gradient of the cost function. This is a two-step process represented by:
\begin{equation}
    \mathbf{w} := \mathbf{w} - \alpha \frac{\partial J(\mathbf{w}, b)}{\partial \mathbf{w}} = \mathbf{w} - \alpha \cdot \nabla_\mathbf{w}J(\mathbf{w}, b)
\end{equation}
\begin{equation}
    b:=b-\alpha \frac{\partial J(\mathbf{w}, b)}{\partial b}
\end{equation}
\subsubsection{Learning Rate}
The parameter $\alpha$ appearing in equations $(5)$ and $(6)$ is called the ``learning rate''. The gradient only tells us the direction in which the cost $J$ increases most steeply, so moving in the opposite direction decreases it; $\alpha$ is what controls how large a step is taken in that direction on each iteration.
\\
Choosing $\alpha$ well matters. If it is too large, each update can overshoot the minimum of $J$, causing the cost to oscillate or even diverge instead of settling down. If it is too small, each step barely moves $\mathbf{w}$ and $b$, and the algorithm needs many more iterations to converge, making training slow. Unlike $\mathbf{w}$ and $b$, $\alpha$ is not learned from the data: it is fixed before training begins, and choosing a suitable value is part of designing the algorithm rather than something gradient descent solves for on its own.

\subsection{Algorithm Structure}

Thus, using equations $(1), (2), (3), (4), (5)$, and $(6)$ for $p$ iterations and $1$ training example, the algorithm can be represented as:\\
Repeat, for $k=1$ to $p$: \{
\begin{align*}
    &\quad z = \mathbf{w}^T\mathbf{x}+b \\
    &\quad a = \hat{y}=\sigma(z) = \frac{1}{1 + e^{-z}} \\
    &\quad J(\mathbf{w},b)= -[y\log a + (1-y)\log{(1-a)}] \\
    &\quad \mathbf{w} = \mathbf{w} - \alpha \frac{\partial J(\mathbf{w}, b)}{\partial \mathbf{w}} = \mathbf{w} - \alpha \cdot \nabla_\mathbf{w}J(\mathbf{w}, b) \\
    &\quad b=b-\alpha \frac{\partial J(\mathbf{w}, b)}{\partial b}=b-\alpha \nabla_{b}J(\mathbf{w},b)
\end{align*}
\}
\\
\\
Note here that repeating the above algorithm multiple times will have no effect as there is only 1 training example, but this idea will be generalized with multiple training examples later.

\subsection{Computation of Derivatives}
To minimize the cost function $J(\mathbf{w},b)$, we must compute its partial derivatives with respect to the two variables $\mathbf{w}$ and $b$. To do so, we can set up a computation graph to calculate the derivatives using the chain rule of calculus.
\subsubsection{Backpropagation}
Backpropagation is the procedure we use to compute these derivatives efficiently. The idea is to read the computation graph from right to left: starting from the cost $J$ at the end, we move backwards one operation at a time, and at each stage we use the chain rule to turn the derivative we have already found into the derivative of the next variable further back. Because every stage reuses the result of the stage before it, we obtain the derivative of the cost with respect to every parameter in a single backward pass, rather than differentiating each parameter from scratch.
\\
\\
Before we begin, it is worth introducing a piece of shorthand that will keep the derivations that follow compact. Every derivative we care about is the derivative of the final quantity, the cost $J$ (which for a single training example is just $\mathcal{L}$), with respect to some intermediate variable. Rather than writing $\frac{\partial J}{\partial z}$, $\frac{\partial J}{\partial \mathbf{w}}$, and $\frac{\partial J}{\partial b}$ out in full every time, we adopt the convention
$$d\text{var} \;\equiv\; \frac{\partial J}{\partial \text{var}},$$
so that $dz = \frac{\partial J}{\partial z}$, $d\mathbf{w} = \frac{\partial J}{\partial \mathbf{w}}$, $db = \frac{\partial J}{\partial b}$, and so on. The prefix $d$ is read as ``the gradient of the cost with respect to'', and it matches the variable names we will later use in code.
%insert computation graph here
$$\mathbf{w} \xrightarrow{} z \xrightarrow[]{} a \xrightarrow{} J$$
We must compute $\frac{\partial J(\mathbf{w},b)}{\partial \mathbf{w}}$ and $\frac{\partial J(\mathbf{w}, b)}{\partial b}$. Using the chain rule, we can write:
\begin{equation}
    \frac{\partial J}{\partial z}=\frac{\partial J}{\partial a} \times \frac{\partial a}{\partial z}
\end{equation}
Similarly,
\begin{equation}
    \frac{\partial J}{\partial \mathbf{w}}=\frac{\partial J}{\partial z} \times \frac{\partial z}{\partial \mathbf{w}}, \quad \frac{\partial J}{\partial b}=\frac{\partial J}{\partial z} \times \frac{\partial z}{\partial b}
\end{equation}
Thus, by combining equations $(7)$ and $(8)$:

\begin{equation}
    \frac{\partial J}{\partial \mathbf{w}}=\frac{\partial J}{\partial a} \times \frac{\partial a}{\partial z} \times \frac{\partial z}{\partial \mathbf{w}}, \quad \frac{\partial J}{\partial b}=\frac{\partial J}{\partial a} \times \frac{\partial a}{\partial z} \times \frac{\partial z}{\partial b}
\end{equation}
Differentiating equation $(3)$ with respect to $a$,
\begin{equation}
    da = \frac{\partial J}{\partial a}=\frac{1-y}{1-a} - \frac{y}{a}
\end{equation}
Differentiating equation $(2)$ with respect to $z$,
\begin{equation}
    \frac{\partial a}{\partial z} = \frac{-1}{(1+e^{-z})^2}\cdot -e^{-z} = \frac{e^{-z}}{(1+e^{-z})^2}
\end{equation}
Using the relation $a = \frac{1}{1 + e^{-z}}$, we can write:
\begin{equation}
    e^{-z} = \frac{1}{a}-1
\end{equation}
Using equations $(11)$ and $(12)$, we can write:
\begin{equation}
    \frac{\partial a}{\partial z} = a(1-a)
\end{equation}
Using equations $(13)$ and $(10)$:
\begin{equation}
    dz = \frac{\partial J}{\partial z}=\left(\frac{1-y}{1-a} - \frac{y}{a}\right)\cdot a(1-a)=a-y
\end{equation}
Additionally, using equation $(1)$ and the fact that $\mathbf{w}=[w_1, w_2, \dots, w_n]$:
$$\frac{\partial z}{\partial w_i}=x_i, \quad \frac{\partial z}{\partial b} = 1$$
Using the above relations and relation $(9)$,
\begin{equation}
    dw_i = \frac{\partial J}{\partial w_i}=(a-y)\cdot x_i, \quad db = \frac{\partial J}{\partial b} = a-y
\end{equation}
Now, these formulae can be used in the iterative process as defined in the Algorithm Structure in Section 3.2.

\subsection{Logistic Regression With Multiple Training Examples}
\subsubsection{Vectorization}

The above model of logistic regression was simplified to process a single input with $n$ features. Now, this model can be generalized using matrices to effectively process $m$ training examples at the same time. This enables the algorithm to perform parallel computations instead of looping over each example one at a time, thus making it more efficient. With $n$ features and $m$ training examples, we can transform $n \times 1$ vector $\mathbf{x}^{(i)}$ (where $i$ represents the $i^{th}$ training example) into an $n \times m$ matrix $\mathbf{X}$, where each column represents a different training example with $n$ features. This can be represented as:
\begin{equation}
    \mathbf{X} = \begin{bmatrix}
\vert & \vert & & \vert \\
\mathbf{x}^{(1)} & \mathbf{x}^{(2)} & \cdots & \mathbf{x}^{(m)} \\
\vert & \vert & & \vert
\end{bmatrix} = \renewcommand{\arraystretch}{1.5}
X = \begin{bmatrix}
x^{(1)}_1 & x^{(2)}_1 & \cdots & x^{(m)}_1 \\
x^{(1)}_2 & x^{(2)}_2 & \cdots & x^{(m)}_2 \\
\vdots & \vdots & \ddots & \vdots \\
x^{(1)}_n & x^{(2)}_n & \cdots & x^{(m)}_n
\end{bmatrix}
\end{equation}
To use matrix $\mathbf{X}$ in equation $(1)$, the scalar $z$ of each training example is collected into a $1 \times m$ row vector $\mathbf{Z}$, and the scalar bias $b$ is added to every one of its entries. This automatic copying of a single value across an entire vector is called ``broadcasting'', and we will rely on it repeatedly. Thus, equation $(1)$ can be rewritten as:
\begin{equation}
    \mathbf{Z} = \mathbf{w}^T\mathbf{X} + b
\end{equation}
Note that $b$ stays a plain scalar here: with only one neuron there is a single bias value, simply broadcast across all $m$ columns of $\mathbf{Z}$. This will change once a layer holds more than one neuron, since each neuron then carries its own bias, and $b$ will become a column vector with one entry per neuron; that vector is in turn broadcast into an effective matrix across the training examples. Throughout, $b$ is written in lowercase (bolded once it becomes a vector) rather than capitalized like $\mathbf{W}$, since it is only ever indexed by neuron and never gains the second, per-input dimension that forces $\mathbf{W}$ to become a full matrix.
\\
Subsequently, $a$ would also turn into a $1 \times m$ row vector $\mathbf{A}$, obtained by applying the sigmoid to each entry of $\mathbf{Z}$ independently:
\begin{equation}
    \mathbf{A}=\sigma(\mathbf{Z})
\end{equation}
where $\mathbf{A}=[a^{(1)}\ a^{(2)}\ \dots\ a^{(m)}]$, and we likewise collect the labels into $\mathbf{Y}=[y^{(1)}\ y^{(2)}\ \dots\ y^{(m)}]$. This would cause the equation for the cost function $J$ to be modified such that it can estimate the average loss in the predictions that it makes.

\begin{equation}
    J=\frac{1}{m} \sum_{i=1}^{m}\mathcal{L}(a^{(i)},y^{(i)}) = -\frac{1}{m} \sum_{i=1}^{m}[y^{(i)}\log a^{(i)} + (1-y^{(i)})\log{(1-a^{(i)})}]
\end{equation}
A subtlety appears the moment we vectorize the backward pass. The sigmoid in $\mathbf{A} = \sigma(\mathbf{Z})$ is applied to each entry separately: the activation $a^{(i)}$ depends on its own $z^{(i)}$ and on no other $z^{(j)}$. As a result, every cross derivative $\frac{\partial a^{(i)}}{\partial z^{(j)}}$ with $i \neq j$ is zero, and the only derivatives that survive are the matching ones, $\frac{\partial a^{(i)}}{\partial z^{(i)}} = \sigma'(z^{(i)})$. When the chain rule carries a gradient from $\mathbf{A}$ back to $\mathbf{Z}$, it would in general sum, for each $z^{(i)}$, the contributions arriving through every $a^{(j)}$; but here all of those cross contributions are zero, so each entry of the incoming gradient is simply scaled by the derivative at that same entry. This entry-by-entry multiplication of two equally shaped vectors (or matrices) is the ``elementwise'', or ``Hadamard'', product, written with the symbol $\odot$ and defined by
$$(\mathbf{u} \odot \mathbf{v})_i = u_i v_i.$$
This is why the elementwise appears as it is the exact form the chain rule takes whenever the function in between acts on each entry independently. Because there are no cross terms left to sum, the usual matrix multiplication reduces to a scaling of matched entries. Writing this out for our case,
$$d\mathbf{Z} = d\mathbf{A} \odot \sigma'(\mathbf{Z}),$$
and substituting the single-example results $\sigma'(z) = a(1-a)$ and $da = \frac{1-y}{1-a} - \frac{y}{a}$ from Section 3.3, the elementwise product collapses neatly to
$$d\mathbf{Z} = \mathbf{A} - \mathbf{Y}.$$
For the sigmoid paired with the cross-entropy loss this cancellation is exact, so the $\odot$ never has to be written out explicitly here; but the very same reasoning, without such a clean cancellation, is exactly what we will need for the hidden layer in Section 4. The gradients of the cost then follow as
$$d\mathbf{w} = \frac{1}{m}\mathbf{X}\,d\mathbf{Z}^T = \frac{1}{m}\mathbf{X}(\mathbf{A}-\mathbf{Y})^T, \qquad db = \frac{1}{m}\sum_{i=1}^{m}(a^{(i)}-y^{(i)}).$$
The update steps for the parameters $\mathbf{w}$ and $b$ then stay the same as before, since their structure is preserved (note that the scalar bias $b$ receives a single averaged gradient over all $m$ examples):
\begin{equation}
    \mathbf{w} := \mathbf{w} - \alpha \frac{\partial J(\mathbf{w}, b)}{\partial \mathbf{w}} = \mathbf{w} - \alpha \cdot d\mathbf{w}
\end{equation}
\begin{equation}
    b:=b-\alpha \frac{\partial J(\mathbf{w}, b)}{\partial b}=b - \alpha \cdot db
\end{equation}

\subsubsection{Algorithm Structure for Multiple Training Examples}

Thus, using equations $(17), (18), (19), (20)$, and $(21)$ for $p$ iterations over all $m$ training examples, the algorithm can be represented as:\\
Repeat, for $k=1$ to $p$: \{
\begin{align*}
    &\quad \mathbf{Z} = \mathbf{w}^T\mathbf{X} + b \\
    &\quad \mathbf{A}=\sigma(\mathbf{Z}) \\
    &\quad J(\mathbf{w},b)= -\frac{1}{m} \sum_{i=1}^{m}[y^{(i)}\log a^{(i)} + (1-y^{(i)})\log{(1-a^{(i)})}] \\
    &\quad d\mathbf{Z} = \mathbf{A}-\mathbf{Y} \\
    &\quad d\mathbf{w}=\frac{1}{m}\mathbf{X}\,d\mathbf{Z}^T = \frac{1}{m}\mathbf{X}(\mathbf{A}-\mathbf{Y})^T \\
    &\quad db=\frac{1}{m}\sum_{i=1}^{m}(a^{(i)}-y^{(i)}) \\
    &\quad \mathbf{w} := \mathbf{w} - \alpha \cdot d\mathbf{w} \\
    &\quad b:=b - \alpha \cdot db
\end{align*}
\}
\\
\\The difference in this implementation compared to that with a single training example is that with multiple training examples, the cost functions, along with the gradients, need to be averaged over the number of training examples taken. This ensures that each individual example contributes to the gradient and the scale of the gradient is also maintained.

\section{Neural Network Structure}
\subsection{Shallow Neural Network}
\subsubsection{Architecture}
A neural network has multiple layers with many "neurons" or nodes inside each layer. Each of these neurons essentially performs logistic regression on its own. A fairly simple structure for a neural network would be one with one hidden layer and one output layer. The structure for this can be built upon by the algorithm that we devised in Section 3.4.2 for multiple training examples.
\\Let the number of layers in the neural network be represented by $L$, and let the components of each layer be denoted by the superscript $^{[l]}$, where $l$ is the layer number. The process for this neural network can be written as:
$$x \xrightarrow{} z^{[1]} \xrightarrow[]{} a^{[1]} \xrightarrow{} z^{[2]} \xrightarrow{} a^{[2]}$$
Let us look at the hidden layer. Let it have $n^{[1]}$ neurons. Thus, these $n^{[1]}$ neurons will all be performing logistic regression, which will be fed into the output layer as $n^{[1]}$ inputs. Since there are $n^{[1]}$ neurons, each neuron must have $n$ parameters $(w_1, w_2, \dots, w_{n})$, one for every input feature. As in Section 3.1, each neuron's weight vector $\mathbf{w}_i^{[1]}$ is stored as an $n \times 1$ column, matching $\mathbf{w}$ in the logistic-regression case; stacking the transposes of the $n^{[1]}$ neurons of the layer as rows gives an $n^{[1]} \times n$ matrix for the whole layer.
$$\mathbf{W}^{[1]} = \begin{bmatrix}
\rule{1cm}{0.4pt} & \mathbf{w}_1^{[1]T} & \rule{1cm}{0.4pt} \\[6pt]
\rule{1cm}{0.4pt} & \mathbf{w}_2^{[1]T} & \rule{1cm}{0.4pt} \\[6pt]
& \vdots & \\[6pt]
\rule{1cm}{0.4pt} & \mathbf{w}_{n^{[1]}}^{[1]T} & \rule{1cm}{0.4pt}
\end{bmatrix} = \begin{bmatrix}
w_{1,1}^{[1]} & w_{1,2}^{[1]} & \cdots & w_{1,n}^{[1]} \\[6pt]
w_{2,1}^{[1]} & w_{2,2}^{[1]} & \cdots & w_{2,n}^{[1]} \\[6pt]
\vdots & \vdots & \ddots & \vdots \\[6pt]
w_{n^{[1]},1}^{[1]} & w_{n^{[1]},2}^{[1]} & \cdots & w_{n^{[1]},n}^{[1]}
\end{bmatrix}$$
Now, if the matrix is substituted in equation $(17)$, this would cause the $1 \times m$ vector to transform into a $n^{[1]} \times m$ matrix $\mathbf{Z}^{[1]}$. As anticipated in Section 3.4.1, since this layer holds $n^{[1]}$ neurons, the bias becomes a column vector $\mathbf{b}^{[1]}$ of shape $n^{[1]} \times 1$, holding one bias per neuron, which is broadcast across all $m$ columns (so that every training example in a given row is offset by the same bias). Thus, the equation can be rewritten as:
\begin{equation}
    \mathbf{Z}^{[1]}=\mathbf{W}^{[1]}\mathbf{X} + \mathbf{b}^{[1]}
\end{equation}
Subsequently, the row vector $\mathbf{A}$ will also be transformed to the matrix $\mathbf{A}^{[1]}$, with the same dimensions as $\mathbf{Z}^{[1]}$, by applying an activation function $g^{[1]}(\cdot)$ to it entry by entry.
\\
At this point the activation must be chosen, and here we make a deliberate departure from logistic regression: we do not reuse the sigmoid for the hidden layer. The sigmoid saturates, meaning that for inputs $z$ that are even moderately large in magnitude its curve flattens out and its derivative $a(1-a)$ becomes vanishingly small. When many such small derivatives are multiplied together through the chain rule during backpropagation, the gradients reaching the earlier layers shrink towards zero and those layers effectively stop learning; this is the well-known ``vanishing gradient'' problem. To avoid it, we use the ``Rectified Linear Unit'' (ReLU) (Nair and Hinton) for the hidden layer, defined as
$$g^{[1]}(z) = \text{ReLU}(z) = \max(0, z).$$
ReLU is trivial to compute, and for every positive input its derivative is exactly $1$, so gradients flow back through the hidden layer undiminished and training stays fast. Its derivative is simply
$$g^{[1]\prime}(z) = \begin{cases} 1 & z > 0 \\ 0 & z < 0 \end{cases}.$$
The hidden-layer activation is therefore
\begin{equation}
    \mathbf{A}^{[1]}=g^{[1]}(\mathbf{Z}^{[1]})=\text{ReLU}(\mathbf{Z}^{[1]})
\end{equation}
Now, $\mathbf{A}^{[1]}$ will be fed into the output layer ($2^{nd}$ layer) as input. Since the output layer has only 1 neuron, the matrix $\mathbf{W}^{[2]}$ will actually be a $1 \times n^{[1]}$ vector, and the bias $b^{[2]}$ is a single scalar broadcast across the $m$ columns.
\begin{equation}
    \mathbf{Z}^{[2]}=\mathbf{W}^{[2]}\mathbf{A}^{[1]} + b^{[2]}
\end{equation}
Subsequently, $\mathbf{Z^{[2]}}$ and $\mathbf{A}^{[2]}$ will take a shape of $1 \times m$. For the output layer, however, we do keep the sigmoid: unlike a hidden unit, the final neuron has to produce a number in $(0,1)$ that we can read as the probability $P(y=1\mid\mathbf{x})$, and the sigmoid is precisely the function that guarantees this. Writing the output activation as $g^{[2]}(z) = \sigma(z)$,
\begin{equation}
    \mathbf{A}^{[2]}=g^{[2]}(\mathbf{Z}^{[2]})=\sigma(\mathbf{Z^{[2]}})
\end{equation}
This division of labour, ReLU in the hidden layer to keep gradients healthy and sigmoid at the output to produce a probability, is the standard choice for a binary classifier.
\subsubsection{Backpropagation}
In this shallow neural network, since there are two layers, we would need to find the derivatives of the parameters of both layers and update them. To do so, we can use computation graphs and backpropagate to find the gradients $d\mathbf{W}^{[1]}$, $d\mathbf{b}^{[1]}$, $d\mathbf{W}^{[2]}$ and $d\mathbf{b}^{[2]}$. The structure of the graph is as follows:
$$\mathbf{W}^{[1]} \xrightarrow{} \mathbf{Z}^{[1]} \xrightarrow[]{} \mathbf{A}^{[1]} \xrightarrow{} \mathbf{W}^{[2]} \xrightarrow{} \mathbf{Z}^{[2]} \xrightarrow[]{} \mathbf{A}^{[2]} \xrightarrow[]{}J$$
The process for finding $d\mathbf{W}^{[2]}$ and $db^{[2]}$ is identical to what we have performed before in the process leading up to equation $(15)$, and the equations just need to be vectorized. Exactly as in the logistic-regression output, the gradient entering the linear step of the output layer is $d\mathbf{Z}^{[2]} = \mathbf{A}^{[2]} - \mathbf{Y}$. Thus, we can write:
\begin{equation}
    d\mathbf{W}^{[2]}=\frac{1}{m}d\mathbf{Z}^{[2]}\mathbf{A}^{[1]T}=\frac{1}{m}(\mathbf{A}^{[2]}-\mathbf{Y})\mathbf{A}^{[1]T}, \quad db^{[2]} =\frac{1}{m}\sum_{i=1}^{m}(a^{(i)}-y^{(i)})
\end{equation}
Now, we can backpropagate the derivatives using the chain rule to find $d\mathbf{W}^{[1]}$ and $d\mathbf{b}^{[1]}$.

\begin{equation}
    \frac{\partial J}{\partial \mathbf{A}^{[1]}}= \frac{\partial J}{\partial \mathbf{Z}^{[2]}} \times\frac{\partial \mathbf{Z}^{[2]}}{\partial \mathbf{A}^{[1]}}
\end{equation}
\begin{equation}
    \frac{\partial J}{\partial \mathbf{Z}^{[1]}}= \frac{\partial J}{\partial \mathbf{A}^{[1]}} \times\frac{\partial \mathbf{A}^{[1]}}{\partial \mathbf{Z}^{[1]}}
\end{equation}
\begin{equation}
    \frac{\partial J}{\partial \mathbf{W}^{[1]}}= \frac{\partial J}{\partial \mathbf{Z}^{[1]}} \times\frac{\partial \mathbf{Z}^{[1]}}{\partial \mathbf{W}^{[1]}}, \quad \frac{\partial J}{\partial \mathbf{b}^{[1]}}= \frac{\partial J}{\partial \mathbf{Z}^{[1]}} \times\frac{\partial \mathbf{Z}^{[1]}}{\partial \mathbf{b}^{[1]}}
\end{equation}
Combining equations $(27), (28)$ and $(29)$, 
\begin{equation}
    \frac{\partial J}{\partial \mathbf{W}^{[1]}}= \frac{\partial J}{\partial \mathbf{Z}^{[2]}} \times\frac{\partial \mathbf{Z}^{[2]}}{\partial \mathbf{A}^{[1]}} \times \frac{\partial \mathbf{A}^{[1]}}{\partial \mathbf{Z}^{[1]}} \times \frac{\partial \mathbf{Z}^{[1]}}{\partial \mathbf{W}^{[1]}}
\end{equation}
\begin{equation}
    \frac{\partial J}{\partial \mathbf{b}^{[1]}}= \frac{\partial J}{\partial \mathbf{Z}^{[2]}} \times\frac{\partial \mathbf{Z}^{[2]}}{\partial \mathbf{A}^{[1]}} \times \frac{\partial \mathbf{A}^{[1]}}{\partial \mathbf{Z}^{[1]}} \times \frac{\partial \mathbf{Z}^{[1]}}{\partial \mathbf{b}^{[1]}}
\end{equation}
We know that:
$$d\mathbf{Z}^{[2]} = \mathbf{A}^{[2]} - \mathbf{Y}$$
Using this in equation $(27)$,
$$\frac{\partial J}{\partial \mathbf{A}^{[1]}} = \left(\frac{\partial \mathbf{Z}^{[2]}}{\partial \mathbf{A}^{[1]}}\right)^T d\mathbf{Z}^{[2]} = \mathbf{W}^{[2]T} d\mathbf{Z}^{[2]}$$
We can also infer that:
$$\frac{\partial \mathbf{A}^{[1]}}{\partial \mathbf{Z}^{[1]}}=g^{[1]'}(\mathbf{Z^{[1]})}$$
Substituting this in equation $(28)$,
$$d\mathbf{Z}^{[1]} = \frac{\partial J}{\partial \mathbf{A}^{[1]}} \odot \frac{\partial \mathbf{A}^{[1]}}{\partial \mathbf{Z}^{[1]}} = \mathbf{W}^{[2]T} d\mathbf{Z}^{[2]} \odot g^{[1]\prime}(\mathbf{Z}^{[1]})$$
This expression shows precisely why one factor is combined as a matrix product and the other as the elementwise product introduced in Section 3.4.1. The first factor, $d\mathbf{A}^{[1]} = \mathbf{W}^{[2]T} d\mathbf{Z}^{[2]}$, comes from the relation $\mathbf{Z}^{[2]} = \mathbf{W}^{[2]}\mathbf{A}^{[1]} + b^{[2]}$. Here each entry of $\mathbf{Z}^{[2]}$ is a weighted sum of every entry of $\mathbf{A}^{[1]}$, so the derivative of $\mathbf{Z}^{[2]}$ with respect to $\mathbf{A}^{[1]}$ has non-zero cross terms, and the chain rule at this step is a full matrix product. The second factor comes from $\mathbf{A}^{[1]} = g^{[1]}(\mathbf{Z}^{[1]})$, where the ReLU acts on each entry of $\mathbf{Z}^{[1]}$ separately. By exactly the argument of Section 3.4.1, its cross derivatives are zero, so this step reduces to the elementwise product with the derivative $g^{[1]\prime}(\mathbf{Z}^{[1]})$. In short, the weight matrix connects every input to every output and therefore requires a matrix product, whereas the activation connects each entry only to itself and therefore requires an elementwise product. For the ReLU, that derivative equals $1$ wherever the corresponding entry of $\mathbf{Z}^{[1]}$ was positive and $0$ otherwise, so the elementwise product keeps the gradient for the units that were active on the forward pass and removes it for the rest. Unlike the output layer, where the sigmoid and the cross-entropy loss cancelled to leave the clean $\mathbf{A}^{[2]} - \mathbf{Y}$, no such cancellation occurs here, so the elementwise product stays in the final expression.
\\
Substituting this in equation $(29)$,
\begin{equation}
    d\mathbf{W}^{[1]} = \frac{1}{m} d\mathbf{Z}^{[1]} \mathbf{X}^T = \frac{1}{m}\left(\mathbf{W}^{[2]T} d\mathbf{Z}^{[2]} \odot g^{[1]\prime}(\mathbf{Z}^{[1]})\right)\mathbf{X}^T
\end{equation}
The bias gradient follows the same substitution into equation $(29)$. Here $\frac{\partial \mathbf{Z}^{[1]}}{\partial \mathbf{b}^{[1]}}=1$, exactly as it was for $b$ in Section 3.3, so $d\mathbf{b}^{[1]}$ is just $d\mathbf{Z}^{[1]}$ itself, averaged across the $m$ training examples. Since $\mathbf{b}^{[1]}$ has one entry per neuron rather than a single scalar, this average is taken across the columns of $d\mathbf{Z}^{[1]}$ only, collapsing the $m$ examples while keeping the $n^{[1]}$ neurons separate:
\begin{equation}
    d\mathbf{b}^{[1]} = \frac{1}{m}\sum_{i=1}^{m} d\mathbf{Z}^{[1](i)}
\end{equation}
where $d\mathbf{Z}^{[1](i)}$ denotes the $i^{\text{th}}$ column of $d\mathbf{Z}^{[1]}$, so that $d\mathbf{b}^{[1]}$ comes out with the same $n^{[1]} \times 1$ shape as $\mathbf{b}^{[1]}$ itself.
\\
With $d\mathbf{W}^{[1]}$, $d\mathbf{b}^{[1]}$, $d\mathbf{W}^{[2]}$, and $d\mathbf{b}^{[2]}$ all computed, both layers are updated using the same gradient descent rule introduced in Section 3.1, applied once per layer:
$$\mathbf{W}^{[l]} := \mathbf{W}^{[l]} - \alpha \cdot d\mathbf{W}^{[l]}, \qquad \mathbf{b}^{[l]} := \mathbf{b}^{[l]} - \alpha \cdot d\mathbf{b}^{[l]}, \qquad l = 1, 2.$$

\subsubsection{Algorithm}
Now, we have all the components necessary to start building the algorithm for this shallow neural network.
Repeat, for $k=1$ to $p$: \{
\begin{align*}
    &\quad \mathbf{Z}^{[1]} = \mathbf{W}^{[1]}\mathbf{X} + \mathbf{b}^{[1]} \\
    &\quad \mathbf{A}^{[1]}=\text{ReLU}(\mathbf{Z}^{[1]}) \\
    &\quad \mathbf{Z}^{[2]} = \mathbf{W}^{[2]}\mathbf{A}^{[1]} + b^{[2]} \\
    &\quad \mathbf{A}^{[2]}=\sigma(\mathbf{Z}^{[2]}) \\
    &\quad J(\mathbf{W}^{[1]},\mathbf{b}^{[1]},\mathbf{W}^{[2]},b^{[2]})= -\frac{1}{m} \sum_{i=1}^{m}[y^{(i)}\log a^{(i)} + (1-y^{(i)})\log{(1-a^{(i)})}] \\
    &\quad d\mathbf{Z}^{[2]} = \mathbf{A}^{[2]}-\mathbf{Y} \\
    &\quad d\mathbf{W}^{[2]}=\frac{1}{m}d\mathbf{Z}^{[2]}\mathbf{A}^{[1]T} \\
    &\quad db^{[2]}=\frac{1}{m}\sum_{i=1}^{m}(a^{(i)}-y^{(i)}) \\
    &\quad d\mathbf{Z}^{[1]} = \mathbf{W}^{[2]T} d\mathbf{Z}^{[2]} \odot g^{[1]\prime}(\mathbf{Z}^{[1]}) \\
    &\quad d\mathbf{W}^{[1]}=\frac{1}{m}d\mathbf{Z}^{[1]}\mathbf{X}^T \\
    &\quad d\mathbf{b}^{[1]}=\frac{1}{m}\sum_{i=1}^{m} d\mathbf{Z}^{[1](i)} \\
    &\quad \mathbf{W}^{[1]} := \mathbf{W}^{[1]} - \alpha \cdot d\mathbf{W}^{[1]} \\
    &\quad \mathbf{b}^{[1]}:=\mathbf{b}^{[1]} - \alpha \cdot d\mathbf{b}^{[1]} \\
    &\quad \mathbf{W}^{[2]} := \mathbf{W}^{[2]} - \alpha \cdot d\mathbf{W}^{[2]} \\
    &\quad b^{[2]}:=b^{[2]} - \alpha \cdot db^{[2]}
\end{align*}
\}
\\
\\
Unlike the earlier algorithms, this one carries two full sets of parameters, $(\mathbf{W}^{[1]}, \mathbf{b}^{[1]})$ and $(\mathbf{W}^{[2]}, b^{[2]})$, that are updated together on every iteration. The same pattern, computing each layer's linear step, applying an activation, and then working backwards one layer at a time to accumulate gradients, extends unchanged to networks with more than two layers; only the number of times it is repeated grows, not the underlying idea.

\subsection{Deep Neural Networks}
We can now generalize the process repeated during the construction of the shallow neural network to form deep neural networks for binary classification. Let the number of layers in the network be $L$, and let the number of nodes in each layer be $n^{[l]}$, where $l$ represents the layer number. Additionally, we can write $\mathbf{X}$ as $\mathbf{A}^{[0]}$ to be able to generalize the equations more effectively.
\\
Thus, for the forward propagation of the network up to the second-to-last layer, we can write:
\\
Repeat, for $l=1$ to $L-1$: \{
\begin{align*}
    &\quad \mathbf{Z}^{[l]}=\mathbf{W}^{[l]} \mathbf{A}^{[l-1]} + \mathbf{b}^{[l]} \\
    &\quad \mathbf{A}^{[l]}=g^{[l]}(\mathbf{Z}^{[l]})
\end{align*}
\}
\\
For the output layer,
$$\mathbf{Z}^{[L]}=\mathbf{W}^{[L]} \mathbf{A}^{[L-1]} + b^{[L]}$$
$$\mathbf{A}^{[L]}=g^{[L]}(\mathbf{Z}^{[L]})$$
As in Section 4.1.1, $g^{[l]}$ for $l < L$ is typically the ReLU, while $g^{[L]}$ is the sigmoid, since the network still ends in a binary classification decision regardless of how many hidden layers precede it.
\\
This shows that the forward propagation steps for the neural network are fairly easy to devise and implement. The backward propagation steps pose a bigger challenge. For the backward propagation, we start at the output layer $L$ and initialize, for each training example $i$,
$$d\mathbf{A}^{[L](i)}=\frac{1}{m}\left[ \frac{1-y^{(i)}}{1-a^{[L](i)}} - \frac{y^{(i)}}{a^{[L](i)}} \right], \quad i = 1, \dots, m,$$
which is simply the single-example result $da = \frac{1-y}{1-a} - \frac{y}{a}$ from equation $(10)$, computed independently in each column and scaled by $\frac{1}{m}$ to match the averaged cost function.
\\
Thereafter, we start the process of computing derivatives backwards for each layer. Since each layer's backward step depends on $d\mathbf{A}^{[l]}$ handed down from the layer after it, this loop must run in the opposite direction to the forward pass, starting at the output layer and working back towards the input:
\\
Repeat, for $l=L$ down to $1$: \{
\begin{align*}
    &\quad d\mathbf{Z}^{[l]}=d\mathbf{A}^{[l]} \odot g^{[l]\prime}(\mathbf{Z}^{[l]}) \\
    &\quad d\mathbf{W}^{[l]} = \frac{1}{m} d\mathbf{Z}^{[l]}\mathbf{A}^{[l-1]T} \\
    &\quad d\mathbf{b}^{[l]}=\frac{1}{m}\sum_{i=1}^{m}d\mathbf{Z}^{[l](i)} \\
    &\quad d\mathbf{A}^{[l-1]}=\mathbf{W}^{[l]T} d\mathbf{Z}^{[l]}
\end{align*}
\}

\section{Implementation}

\subsection{Familiarization With the Dataset}
The Higgs Boson dataset has $200,000$ rows with $30$ features of events, along with a binary label of whether an event was a Higgs Boson signal event $(1)$ or background noise $(0)$.
\\
At the LHC (Large Hadron Collider), protons are smashed together at extremely high energy, producing sprays of particles that detectors measure. Most of these collisions produce ordinary, well-understood physics processes ("background"). Very rarely, a collision produces a Higgs boson, which almost immediately decays into other particles before it can be directly observed. The pattern of particles that come out afterward (their energies, momenta, angles) could indicate whether the particles come from a Higgs decay versus other common physics processes.
\\
The split of signal events to background noise in the dataset is $\sim 0.657: 0.343$, i.e. always predicting "background" would yield an accuracy of $65.7\%$.

\subsection{Pre-processing the Data}
Before the data can be fed into any model, it first needs to be pre-processed to ensure that it is in the correct format for the required input. To do this, some features of the data, such as the "EventID" and "Weight", were removed from the data as they would not have had any impact while training. Additionally, the data was converted from a \verb|.csv| format to a \verb|NumPy| array to aid in computation. 
\\\\
Further, the data was also normalized. This was done for a few reasons. Firstly, since different features are spread between different orders of magnitude, those that are large in raw value would dominate the $\mathbf{w}^T\mathbf{x} + b$ computation. Thus, the magnitude of the weight given to a feature would represent how it impacts the final outcome rather than how big or small the feature is. This also prevents saturation in activation functions, especially in the sigmoid activation, as extreme values would push the function to either $0$ or $1$, which has a very low rate of change, and there are greater chances for the weights to be "stuck" in such low-gradient areas.
\\\\
An $80:20$ split $(200,000:50,000)$ was used to separate the $250,000$ rows of data into training and testing sets, according to common convention.

\subsection{Initialization of Parameters}
Since initializing the weight matrices to zero would essentially mean that every neuron in the layer would give an identical result, the values of the weight matrices of each layer were initialized to random values using the \verb|np.random.randn()| function and were scaled by a factor of $\sqrt{\frac{2}{l}}$, where $l$ is the number of neurons in the previous layer of the neural network. This is called the ``He'' initialization, introduced by Kaiming He and his colleagues in 2015 specifically to work with ReLU activations (He et al.). This type of scaling helps prevent vanishing or exploding gradients.


\subsection{Testing}
\subsubsection{Logistic Regression on a Single Neuron}
To put the algorithm to the test, a logistic regression model was implemented, on which the data was trained. This model showed an accuracy of $\sim 75.0\% $ on the training set and $\sim 74.6\%$ accuracy on the test set at a learning rate of $r=0.01$ for $10,000$ iterations, which is already better than the baseline accuracy of $65.7\%$, indicating that a refined model could possibly yield better predictions. With these results, I decided to implement a neural network to understand how far accurate the predictions could be.

\subsection{Deep Neural Network}
Looking at the improvements from the single neuron model, I decided to implement a deep neural network to test its accuracy. It was composed of 5 layers, and the following architecture of neurons in each layer:
$$\mathbf{A}^{[0]} \xrightarrow{} 20 \xrightarrow{} 10 \xrightarrow{} 7 \xrightarrow{} 4 \xrightarrow{} 1$$

\subsection{Hyperparamater Tuning}
\subsubsection{Learning Rate}
Initially, when I tried training it at a learning rate of $r=0.01$, I found the accuracy to be $\sim 78\%$ for both the training and the test sets. Then, looking at the cost curve, I realized that the model was far away from convergence. Thus, I decided to train the model on a number of learning rates:
$$[0.01,\space 0.025, \space  0.075,\space 0.1,\space 0.2,\space 0.3,\space 0.4,\space 0.5,\space 0.6,\space 0.7\space]$$
The cost curves of the models and their accuracies can be seen in the following figures.
%insert cost curves and accuracies

\begin{figure}[htbp]
    \centering % Centers the image on the page
    \includegraphics[width=0.9\textwidth]{output.png} % Adjust size and insert file
    \caption{Cost curve for all the different learning rates tested} % Add a caption
    \label{fig:cost_curve_lrs} % Add a unique label for cross-referencing
\end{figure} 

Figure 1 shows that a higher learning rate converges faster and generally at a lower cost for a fixed number of iterations. We do have to note that this cannot be generalized as once a learning rate gets too big, it can fail to find the optima as it could take steps that would be too big as well. When that is the case, the graph for cost would not show a smooth curve but one that has many spikes in between, and early signs of these spikings can be seen if one looks closely at the plot for $r=0.7$. 

\begin{figure}[htbp]
    \centering % Centers the image on the page
    \includegraphics[width=0.9\textwidth]{full_gd_epochs.png} % Adjust size and insert file
    \caption{Graph of the train and test accuracies for each learning rate} % Add a caption
    \label{fig:cost_curve_lrs} % Add a unique label for cross-referencing
\end{figure}

Looking at Figure 2, the model with a learning rate of $r=0.6$ showed the highest training accuracy of $\sim 84.3\%$, along with a test accuracy of $\sim 83.9\%$. This echoes the findings from Figure 1, as we see that a higher learning rate has enabled the model to reach higher accuracy within a fixed number of iterations but the highest learning rate is not necessarily the best one to use. 
\\\\
This showed me that finding the best neural network model for a specific dataset and use case is not extremely obvious, and requires quite a lot of trial and error.

\subsubsection{Model Architecture}
Since experimenting with the learning rate yielded a variety of results, I thought I could experiment with different model architectures as well to explore how it can affect and help improve accuracy. 
\\\\
There were some changes that were made to the implementation in code with this decision. Previously, training these neural networks for 11 learning rates was taking a very long time (around 5 hours) on my laptop. Examining this, I knew there had to be a method to train faster as I had seen frameworks like TensorFlow train such models with a similar dataset size in a much shorter time. While researching, I came across a python package called \verb|cupy| (CuPy Developers), which has all the same functions as \verb|numpy|, along with the ability to compute matrix multiplications using the GPU. This also needed no changes to the codebase of my neural network, as I could simply swap packages while preserving all functions and code I had written previously. Thus, I was able to easily implement this change and train on my modest laptop GPU.
\\\\
Additionally, I also came across the idea of training the network in mini-batches while researching. In this method, the data is divided into a number of small batches, and the model processes each batch alone and updates parameters after the pass is run for each batch. An epoch is completed after each batch is processed, and the model is trained for a number of epochs similar to the number of iterations we used earlier). This is essentially the process frameworks like TensorFlow and PyTorch use (TensorFlow Developers), which enables them to train very large networks in a relatively short amount of time. This is computationally cheaper, as the model does not have to perform a full pass with the whole $200,000$ rows of the dataset. Instead, it gets the opportunity to update its parameters more frequently, which in many cases results in a faster path to convergence. Additionally, since it does not process the full size of the dataset at once, it is more memory efficient and lets models train on a small amount of available memory as well.
\\\\
Integrating both of these changes helped me train 8 different model architectures that were more complex in a much shorter amount of time ($\sim 75$ minutes). The architectures are seen in the following table.

\begin{table}[h]
\centering
\begin{tabular}{l l c r}
\hline
\textbf{Architecture} & \textbf{Hidden layer widths} & \textbf{Total layers} & \textbf{Parameters} \\
\hline
Baseline (previous)     & 20, 10, 7, 4          & 5 & 944 \\
Double width             & 40, 20, 14, 8         & 5 & 2,483 \\
Extra layer               & 40, 20, 10, 7, 4      & 6 & 2,384 \\
Two extra layers          & 80, 40, 20, 10, 7, 4  & 7 & 6,864 \\
Wide                      & 128, 64                & 3 & 12,289 \\
Very wide                 & 256, 128               & 3 & 40,961 \\
Deep and wide              & 128, 64, 32, 16, 8    & 6 & 14,977 \\
Oversized                  & 512, 256, 128, 64     & 5 & 188,417 \\
\hline
\end{tabular}
\caption{Architectures tested, given by the number of neurons in each hidden layer, along with total parameter count. Parameter count for layer $l$ is $n^{[l-1]} \times n^{[l]} + n^{[l]}$ (weights plus biases), summed across all layers. All networks share the same 30-feature input layer and single-neuron sigmoid output layer.}
\label{tab:architectures}
\end{table}
As I started training these models on the best learning rate I had before $r=0.6$, I started getting \verb|nan| in the cost computation I was logging, which was a signal that the cost had exploded. Upon researching, I realized that models with mini-batching implemented had to be trained on far lower learning rates, and complex models needed to have more batches to prevent the cost from exploding. Thus, I chose to train my models at a learning rate of $r=0.01$, for $400$ epochs by dividing the data into batches of $1024$. The cost curves and accuracies of the models are shown in the following figures. An important point to note here is that I was not able to train my "oversized" model fully with this learning rate, as the cost exploded.

\begin{figure}[htbp]
    \centering % Centers the image on the page
    \includegraphics[width=0.9\textwidth]{cost_with_epochs_architecture.png} % Adjust size and insert file
    \caption{Cost curve for all the different model architectures tested} % Add a caption
    \label{fig:cost_curve_lrs} % Add a unique label for cross-referencing
\end{figure}

\begin{figure}[htbp]
    \centering % Centers the image on the page
    \includegraphics[width=0.9\textwidth]{final_acc_graph.png} % Adjust size and insert file
    \caption{Train and test accuracies of the different model architectures tested} % Add a caption
    \label{fig:cost_curve_lrs} % Add a unique label for cross-referencing
\end{figure}

Upon analyzing the graphs, one can notice that the best result was obtained from the "deep and wide" model, with $6$ layers and $14,977$ parameters, which yielded a training accuracy of $85.4\%$ but did not score too well in terms of the testing accuracy $(83.4\%)$. There is a clear observation I could make from this experiment. As the model architecture gets deeper and wider, the training accuracy continues to grow in minor increments, while the testing accuracy stays at a plateau. This is also evidenced by the cost curve of all of the models, especially the "oversized" one, which could reach to much lower costs while training before the cost exploded. Thus, it can be concluded that increasing the number of parameters has an effect on minimizing the cost more efficiently.

\subsection{Evaluation}
Since the dataset was used for the ATLAS Challenge competition, there have been many attempts at training different types of machine learning algorithms to try to achieve the best accuracy at prediction. I found a paper by S. Raza Ahmad which similarly implemented a deep neural network from scratch using NumPy for this same competition, achieving an accuracy of $83\%$ using a 4-hidden-layer network with 300 units per layer, sigmoid activations throughout, and layer-wise pretraining using stacked autoencoders, trained across a 12-node distributed cluster over 7-10 days (Ahmad). It is worth noting that this comparison is not perfectly like-for-like given the substantial differences in architecture, compute, and training methodology; nonetheless, my comparatively lightweight, single-model network achieved a marginally higher accuracy of $\sim 84\%$.
\\\\
This suggests that, for this dataset, a comparatively small, single from-scratch network trained on a laptop GPU was able to match, or even marginally exceed, the accuracy of a much larger and more heavily engineered deep learning solution built with far greater compute and time. This does not mean my implementation is inherently better; the referenced work explored a broader set of optimization techniques (momentum, RMSProp, learning rate annealing, layer-wise pretraining) that were outside the scope of this project. Rather, it suggests that for this particular dataset and feature set, the primary bottleneck on achievable accuracy is likely the amount of learnable signal present in the 30 features themselves, rather than model scale or training sophistication.

\section{Conclusion}

Thus, from this project, I have learned how to implement and fine-tune a neural network from scratch. Understanding the underlying mathematical mechanics of neural networks helped me understand how to construct and structure one, from making a choice of the activation functions to implementing it all in code. This gave me a much deeper intuition about how neural networks work as opposed to stacking some layers with a package like TensorFlow.
\\\\
Deriving backpropagation by hand, from the single-neuron case in to the fully vectorized version, helped me understand concepts that previously felt like black-box function calls. Additionally, testing the model empirically taught me just as much as the derivations did. Comparing learning rates and model architecture showed me that convergence speed and the quality of the final solution are two different problems that might not be solved together. 
\\\\
Going forward, I plan to apply this same from-scratch approach to convolutional neural networks, and learn about model assessment, evaluation, and regularization to build a full machine learning project.

\clearpage
\begin{center}
Works Cited
\end{center}

\noindent\hangindent=0.5in\hangafter=1
Ahmad, S. Raza. ``Technical Report of Participation in Higgs Boson Machine Learning Challenge.'' \textit{arXiv}, 9 Oct. 2015, \url{arxiv.org/abs/1510.02674}.
\par\bigskip

\noindent\hangindent=0.5in\hangafter=1
CuPy Developers. ``CuPy: NumPy \& SciPy for GPU.'' \textit{CuPy}, \url{cupy.dev}.
\par\bigskip

\noindent\hangindent=0.5in\hangafter=1
He, Kaiming, et al. ``Delving Deep into Rectifiers: Surpassing Human-Level Performance on ImageNet Classification.'' \textit{arXiv}, 6 Feb. 2015, \url{arxiv.org/abs/1502.01852}.
\par\bigskip

\noindent\hangindent=0.5in\hangafter=1
``Higgs Boson Machine Learning Challenge.'' \textit{Kaggle}, 2014, \url{www.kaggle.com/c/higgs-boson}.
\par\bigskip

\noindent\hangindent=0.5in\hangafter=1
Nair, Vinod, and Geoffrey E. Hinton. ``Rectified Linear Units Improve Restricted Boltzmann Machines.'' \textit{Proceedings of the 27th International Conference on Machine Learning}, 2010, \url{www.cs.toronto.edu/~fritz/absps/reluICML.pdf}.
\par\bigskip

\noindent\hangindent=0.5in\hangafter=1
Ng, Andrew. ``Neural Networks and Deep Learning.'' \textit{Coursera}, DeepLearning.AI, \url{www.coursera.org/learn/neural-networks-deep-learning}.
\par\bigskip

\noindent\hangindent=0.5in\hangafter=1
TensorFlow Developers. ``TensorFlow Core.'' \textit{TensorFlow}, \url{www.tensorflow.org}.

\end{document}
